% Options for packages loaded elsewhere
\PassOptionsToPackage{unicode}{hyperref}
\PassOptionsToPackage{hyphens}{url}
%
\documentclass[
  11pt,
]{article}
\usepackage{amsmath,amssymb}
\usepackage{iftex}
\ifPDFTeX
  \usepackage[T1]{fontenc}
  \usepackage[utf8]{inputenc}
  \usepackage{textcomp} % provide euro and other symbols
\else % if luatex or xetex
  \usepackage{unicode-math} % this also loads fontspec
  \defaultfontfeatures{Scale=MatchLowercase}
  \defaultfontfeatures[\rmfamily]{Ligatures=TeX,Scale=1}
\fi
\usepackage{lmodern}
\ifPDFTeX\else
  % xetex/luatex font selection
\fi
% Use upquote if available, for straight quotes in verbatim environments
\IfFileExists{upquote.sty}{\usepackage{upquote}}{}
\IfFileExists{microtype.sty}{% use microtype if available
  \usepackage[]{microtype}
  \UseMicrotypeSet[protrusion]{basicmath} % disable protrusion for tt fonts
}{}
\makeatletter
\@ifundefined{KOMAClassName}{% if non-KOMA class
  \IfFileExists{parskip.sty}{%
    \usepackage{parskip}
  }{% else
    \setlength{\parindent}{0pt}
    \setlength{\parskip}{6pt plus 2pt minus 1pt}}
}{% if KOMA class
  \KOMAoptions{parskip=half}}
\makeatother
\usepackage{xcolor}
\usepackage[margin=2.2cm]{geometry}
\usepackage{longtable,booktabs,array}
\usepackage{calc} % for calculating minipage widths
% Correct order of tables after \paragraph or \subparagraph
\usepackage{etoolbox}
\makeatletter
\patchcmd\longtable{\par}{\if@noskipsec\mbox{}\fi\par}{}{}
\makeatother
% Allow footnotes in longtable head/foot
\IfFileExists{footnotehyper.sty}{\usepackage{footnotehyper}}{\usepackage{footnote}}
\makesavenoteenv{longtable}
\setlength{\emergencystretch}{3em} % prevent overfull lines
\providecommand{\tightlist}{%
  \setlength{\itemsep}{0pt}\setlength{\parskip}{0pt}}
\setcounter{secnumdepth}{-\maxdimen} % remove section numbering
\ifLuaTeX
  \usepackage{selnolig}  % disable illegal ligatures
\fi
\IfFileExists{bookmark.sty}{\usepackage{bookmark}}{\usepackage{hyperref}}
\IfFileExists{xurl.sty}{\usepackage{xurl}}{} % add URL line breaks if available
\urlstyle{same}
\hypersetup{
  pdftitle={Worksheet 3B: Quadratics --- Factoring and Solving},
  hidelinks,
  pdfcreator={LaTeX via pandoc}}

\title{Worksheet 3B: Quadratics --- Factoring and Solving}
\author{}
\date{}

\begin{document}
\maketitle

\hypertarget{part-a-factoring-the-reverse-of-foil-warm-up}{%
\subsection{Part A: Factoring --- The Reverse of FOIL
(Warm-up)}\label{part-a-factoring-the-reverse-of-foil-warm-up}}

In Worksheet 3A, you expanded \((t + 3)(t + 4)\) to get
\(t^2 + 7t + 12\).

Now we go backwards: given \(t^2 + 7t + 12\), find the two binomials.

\textbf{The Key Insight:}

To factor \(x^2 + bx + c\), find two numbers that: - \textbf{Multiply}
to give \(c\) (the constant term) - \textbf{Add} to give \(b\) (the
coefficient of \(x\))

\textbf{Example:} Factor \(x^2 + 7x + 12\)

\begin{itemize}
\tightlist
\item
  Need two numbers that multiply to 12 and add to 7
\item
  Factors of 12: (1,12), (2,6), (3,4)
\item
  Which pair adds to 7? → 3 and 4
\item
  Answer: \((x + 3)(x + 4)\)
\end{itemize}

\textbf{Factor each expression. Check your answer by FOILing!}

\textbf{1.} \(t^2 + 5t + 6\)

\begin{itemize}
\tightlist
\item
  Multiply to \_\_\_, add to \_\_\_
\item
  Numbers: \_\_\_ and \_\_\_
\item
  Factored: (\_\_\_)(\_\_\_)
\end{itemize}

\textbf{2.} \(p^2 + 8p + 15\)

\textbf{3.} \(m^2 + 9m + 20\)

\textbf{4.} \(r^2 + 11r + 24\)

\textbf{5.} \(k^2 + 7k + 10\)

\textbf{6.} \(n^2 + 12n + 35\)

\begin{center}\rule{0.5\linewidth}{0.5pt}\end{center}

\hypertarget{part-b-factoring-with-negative-middle-terms-8-problems}{%
\subsection{Part B: Factoring with Negative Middle Terms (8
problems)}\label{part-b-factoring-with-negative-middle-terms-8-problems}}

When the middle term is negative but the constant is positive, both
numbers are negative.

\textbf{Example:} Factor \(x^2 - 7x + 12\)

\begin{itemize}
\tightlist
\item
  Need two numbers that multiply to +12 and add to -7
\item
  Both must be negative (negative × negative = positive)
\item
  Numbers: -3 and -4
\item
  Answer: \((x - 3)(x - 4)\)
\end{itemize}

\textbf{Factor completely:}

\textbf{7.} \(t^2 - 5t + 6\)

\textbf{8.} \(p^2 - 9p + 14\)

\textbf{9.} \(m^2 - 8m + 15\)

\textbf{10.} \(r^2 - 10r + 21\)

\textbf{11.} \(k^2 - 11k + 28\)

\textbf{12.} \(n^2 - 13n + 36\)

\textbf{13.} \(w^2 - 6w + 9\) --- \emph{What's special about this one?}

\textbf{14.} \(s^2 - 14s + 49\) --- \emph{Another special one!}

\begin{center}\rule{0.5\linewidth}{0.5pt}\end{center}

\hypertarget{part-c-factoring-with-negative-constants-10-problems}{%
\subsection{Part C: Factoring with Negative Constants (10
problems)}\label{part-c-factoring-with-negative-constants-10-problems}}

When the constant term is negative, one number is positive and one is
negative.

\textbf{Example:} Factor \(x^2 + 2x - 15\)

\begin{itemize}
\tightlist
\item
  Need two numbers that multiply to -15 and add to +2
\item
  One positive, one negative
\item
  Factors of 15: (1,15), (3,5)
\item
  Which pair gives +2 when combined? → +5 and -3
\item
  Answer: \((x + 5)(x - 3)\)
\end{itemize}

\textbf{Factor completely:}

\textbf{15.} \(t^2 + 3t - 10\)

\textbf{16.} \(p^2 + 2p - 8\)

\textbf{17.} \(m^2 - 2m - 15\)

\textbf{18.} \(r^2 - 3r - 18\)

\textbf{19.} \(k^2 + 4k - 21\)

\textbf{20.} \(n^2 - 5n - 24\)

\textbf{21.} \(w^2 + w - 30\)

\textbf{22.} \(s^2 - 6s - 27\)

\textbf{23.} \(h^2 + 7h - 44\)

\textbf{24.} \(u^2 - u - 56\)

\begin{center}\rule{0.5\linewidth}{0.5pt}\end{center}

\hypertarget{part-d-special-factoring-patterns-10-problems}{%
\subsection{Part D: Special Factoring Patterns (10
problems)}\label{part-d-special-factoring-patterns-10-problems}}

\hypertarget{pattern-1-difference-of-squares}{%
\subsubsection{Pattern 1: Difference of
Squares}\label{pattern-1-difference-of-squares}}

\[a^2 - b^2 = (a + b)(a - b)\]

\emph{Look for two perfect squares separated by subtraction.}

\textbf{Factor:}

\textbf{25.} \(t^2 - 25\)

\textbf{26.} \(p^2 - 81\)

\textbf{27.} \(m^2 - 100\)

\textbf{28.} \(4r^2 - 9\)

\textbf{29.} \(16k^2 - 49\)

\textbf{30.} \(25n^2 - 1\)

\hypertarget{pattern-2-perfect-square-trinomials}{%
\subsubsection{Pattern 2: Perfect Square
Trinomials}\label{pattern-2-perfect-square-trinomials}}

\[a^2 + 2ab + b^2 = (a + b)^2\] \[a^2 - 2ab + b^2 = (a - b)^2\]

\emph{Check: Is the middle term exactly twice the product of the square
roots of the first and last terms?}

\textbf{Factor:}

\textbf{31.} \(w^2 + 10w + 25\)

\textbf{32.} \(s^2 - 12s + 36\)

\textbf{33.} \(h^2 + 16h + 64\)

\textbf{34.} \(4u^2 + 12u + 9\)

\begin{center}\rule{0.5\linewidth}{0.5pt}\end{center}

\hypertarget{part-e-factoring-when-the-leading-coefficient-1-12-problems}{%
\subsection{Part E: Factoring When the Leading Coefficient ≠ 1 (12
problems)}\label{part-e-factoring-when-the-leading-coefficient-1-12-problems}}

When you have \(ax^2 + bx + c\) with \(a \neq 1\), factoring is
trickier.

\textbf{Method: AC Method (Factor by Grouping)}

\textbf{Example:} Factor \(2x^2 + 7x + 3\)

\begin{enumerate}
\def\labelenumi{\arabic{enumi}.}
\tightlist
\item
  Multiply \(a \cdot c\): \(2 \cdot 3 = 6\)
\item
  Find two numbers that multiply to 6 and add to 7: → 6 and 1
\item
  Rewrite the middle term: \(2x^2 + 6x + 1x + 3\)
\item
  Group and factor: \(2x(x + 3) + 1(x + 3)\)
\item
  Factor out common binomial: \((2x + 1)(x + 3)\)
\end{enumerate}

\textbf{Factor completely using the AC method:}

\textbf{35.} \(2t^2 + 5t + 3\)

\begin{itemize}
\tightlist
\item
  \(a \cdot c ={}\) \_\_\_
\item
  Numbers that multiply to \_\_\_ and add to 5: \_\_\_ and \_\_\_
\item
  Rewrite: \(2t^2 + \_\_\_t + \_\_\_t + 3\)
\item
  Factor: (\_\_\_)(\_\_\_)
\end{itemize}

\textbf{36.} \(3p^2 + 10p + 8\)

\textbf{37.} \(2m^2 + 9m + 4\)

\textbf{38.} \(5r^2 + 13r + 6\)

\textbf{39.} \(3k^2 - 7k + 2\)

\textbf{40.} \(2n^2 - 11n + 5\)

\textbf{41.} \(4w^2 + 4w - 3\)

\textbf{42.} \(6s^2 - 7s - 3\)

\textbf{43.} \(3h^2 - 10h - 8\)

\textbf{44.} \(4u^2 - 25\) --- \emph{Difference of squares!}

\textbf{45.} \(9v^2 + 12v + 4\) --- \emph{Perfect square trinomial!}

\textbf{46.} \(2q^2 - 18\) --- \emph{Factor out GCF first!}

\begin{center}\rule{0.5\linewidth}{0.5pt}\end{center}

\hypertarget{part-f-always-look-for-gcf-first-6-problems}{%
\subsection{Part F: Always Look for GCF First! (6
problems)}\label{part-f-always-look-for-gcf-first-6-problems}}

Before using any other technique, always check for a \textbf{Greatest
Common Factor}.

\textbf{Example:} Factor \(3x^2 + 12x + 9\)

\begin{enumerate}
\def\labelenumi{\arabic{enumi}.}
\tightlist
\item
  GCF of 3, 12, 9 is 3
\item
  Factor out 3: \(3(x^2 + 4x + 3)\)
\item
  Factor the trinomial: \(3(x + 1)(x + 3)\)
\end{enumerate}

\textbf{Factor completely (GCF first, then continue):}

\textbf{47.} \(2t^2 + 10t + 12\)

\textbf{48.} \(5p^2 - 20p + 15\)

\textbf{49.} \(3m^2 - 27\)

\textbf{50.} \(4r^3 + 8r^2 - 12r\)

\textbf{51.} \(6k^2 - 24k + 24\)

\textbf{52.} \(2n^3 - 8n^2 - 10n\)

\begin{center}\rule{0.5\linewidth}{0.5pt}\end{center}

\hypertarget{part-g-solving-by-factoring-zero-product-property-12-problems}{%
\subsection{Part G: Solving by Factoring --- Zero Product Property (12
problems)}\label{part-g-solving-by-factoring-zero-product-property-12-problems}}

\textbf{The Zero Product Property:}

If \(A \cdot B = 0\), then \(A = 0\) or \(B = 0\) (or both).

This lets us solve factored equations!

\textbf{Example:} Solve \(x^2 - 5x + 6 = 0\)

\begin{enumerate}
\def\labelenumi{\arabic{enumi}.}
\tightlist
\item
  Factor: \((x - 2)(x - 3) = 0\)
\item
  Apply Zero Product Property: \(x - 2 = 0\) or \(x - 3 = 0\)
\item
  Solve each: \(x = 2\) or \(x = 3\)
\end{enumerate}

\textbf{Solve by factoring:}

\textbf{53.} \(t^2 + 7t + 12 = 0\)

\textbf{54.} \(p^2 - 6p + 8 = 0\)

\textbf{55.} \(m^2 + 2m - 15 = 0\)

\textbf{56.} \(r^2 - 9 = 0\)

\textbf{57.} \(k^2 - 10k + 25 = 0\) --- \emph{How many solutions?}

\textbf{58.} \(2n^2 + 5n - 3 = 0\)

\textbf{59.} \(3w^2 - 7w + 2 = 0\)

\textbf{60.} \(s^2 + 8s = 0\) --- \emph{Don't forget: s is a common
factor!}

\textbf{61.} \(h^2 = 16\) --- \emph{Rearrange first!}

\textbf{62.} \(u^2 - 3u = 10\) --- \emph{Rearrange to standard form
first!}

\textbf{63.} \(2v^2 + 6v = -4\) --- \emph{Rearrange, then factor out
GCF}

\textbf{64.} \(t(t - 4) = 12\) --- \emph{Expand first, then rearrange}

\begin{center}\rule{0.5\linewidth}{0.5pt}\end{center}

\hypertarget{part-h-completing-the-square-10-problems}{%
\subsection{Part H: Completing the Square (10
problems)}\label{part-h-completing-the-square-10-problems}}

Sometimes factoring doesn't work nicely. \textbf{Completing the square}
always works.

\textbf{The Idea:} Turn \(x^2 + bx\) into a perfect square by adding
\(\left(\frac{b}{2}\right)^2\)

\textbf{Example:} Solve \(x^2 + 6x = 7\)

\begin{enumerate}
\def\labelenumi{\arabic{enumi}.}
\tightlist
\item
  Take half of 6, square it: \(\left(\frac{6}{2}\right)^2 = 9\)
\item
  Add 9 to both sides: \(x^2 + 6x + 9 = 7 + 9\)
\item
  Factor left side: \((x + 3)^2 = 16\)
\item
  Take square root: \(x + 3 = \pm 4\)
\item
  Solve: \(x = -3 + 4 = 1\) or \(x = -3 - 4 = -7\)
\end{enumerate}

\textbf{Solve by completing the square:}

\textbf{65.} \(t^2 + 4t = 5\)

\begin{itemize}
\tightlist
\item
  Half of 4, squared: \_\_\_
\item
  Add to both sides: \(t^2 + 4t + \_\_\_ = 5 + \_\_\_\)
\item
  Factor: \((t + \_\_\_)^2 = \_\_\_\)
\item
  Solutions: \(t ={}\) \_\_\_
\end{itemize}

\textbf{66.} \(p^2 - 8p = 9\)

\textbf{67.} \(m^2 + 10m = -21\)

\textbf{68.} \(r^2 - 2r = 15\)

\textbf{69.} \(k^2 + 6k + 5 = 0\) --- \emph{Move 5 to the right first}

\textbf{70.} \(n^2 - 4n - 12 = 0\)

\textbf{71.} \(w^2 + 3w = 4\) --- \emph{Yes, you can complete the square
with odd numbers!}

\textbf{72.} \(s^2 - 5s = 6\)

\textbf{73.} \(h^2 + 8h + 20 = 0\) --- \emph{What happens here?}

\textbf{74.} \(u^2 - 6u + 10 = 0\) --- \emph{And here?}

\begin{center}\rule{0.5\linewidth}{0.5pt}\end{center}

\hypertarget{part-i-the-quadratic-formula-10-problems}{%
\subsection{Part I: The Quadratic Formula (10
problems)}\label{part-i-the-quadratic-formula-10-problems}}

For any quadratic \(ax^2 + bx + c = 0\):

\[\boxed{x = \frac{-b \pm \sqrt{b^2 - 4ac}}{2a}}\]

\textbf{Example:} Solve \(2x^2 + 5x - 3 = 0\)

\begin{itemize}
\tightlist
\item
  \(a = 2\), \(b = 5\), \(c = -3\)
\item
  \(x = \frac{-5 \pm \sqrt{25 - 4(2)(-3)}}{2(2)} = \frac{-5 \pm \sqrt{25 + 24}}{4} = \frac{-5 \pm \sqrt{49}}{4} = \frac{-5 \pm 7}{4}\)
\item
  \(x = \frac{-5 + 7}{4} = \frac{1}{2}\) or
  \(x = \frac{-5 - 7}{4} = -3\)
\end{itemize}

\textbf{Solve using the quadratic formula:}

\textbf{75.} \(t^2 - 5t + 6 = 0\) --- \emph{Verify by factoring too!}

\textbf{76.} \(p^2 + 3p - 10 = 0\)

\textbf{77.} \(2m^2 + 7m + 3 = 0\)

\textbf{78.} \(3r^2 - 2r - 1 = 0\)

\textbf{79.} \(k^2 - 4k + 1 = 0\) --- \emph{Answer won't be a nice
integer}

\textbf{80.} \(n^2 + 2n - 2 = 0\)

\textbf{81.} \(w^2 + 4w + 5 = 0\) --- \emph{What happens under the
square root?}

\textbf{82.} \(2s^2 - 3s + 5 = 0\) --- \emph{Same question!}

\textbf{83.} \(h^2 - 6h + 9 = 0\) --- \emph{What's special about this
one?}

\textbf{84.} \(4u^2 + 4u + 1 = 0\)

\begin{center}\rule{0.5\linewidth}{0.5pt}\end{center}

\hypertarget{part-j-the-discriminant-predicting-solutions-8-problems}{%
\subsection{Part J: The Discriminant --- Predicting Solutions (8
problems)}\label{part-j-the-discriminant-predicting-solutions-8-problems}}

The expression under the square root, \(b^2 - 4ac\), is called the
\textbf{discriminant}.

It tells you what kind of solutions to expect:

\begin{longtable}[]{@{}ll@{}}
\toprule\noalign{}
Discriminant & Number of Real Solutions \\
\midrule\noalign{}
\endhead
\bottomrule\noalign{}
\endlastfoot
\(b^2 - 4ac > 0\) & Two different real solutions \\
\(b^2 - 4ac = 0\) & One real solution (repeated) \\
\(b^2 - 4ac < 0\) & No real solutions (complex numbers) \\
\end{longtable}

\textbf{Without solving, calculate the discriminant and predict the
number of real solutions:}

\textbf{85.} \(t^2 + 5t + 6 = 0\)

\begin{itemize}
\tightlist
\item
  \(b^2 - 4ac = 25 - 24 = 1\)
\item
  Discriminant is \_\_\_\_\_\_\_\_ (positive/zero/negative)
\item
  Number of real solutions: \_\_\_
\end{itemize}

\textbf{86.} \(p^2 - 8p + 16 = 0\)

\textbf{87.} \(m^2 + 2m + 5 = 0\)

\textbf{88.} \(r^2 - 3r - 10 = 0\)

\textbf{89.} \(k^2 + k + 1 = 0\)

\textbf{90.} \(3n^2 - 2n - 1 = 0\)

\textbf{91.} \(4w^2 - 12w + 9 = 0\)

\textbf{92.} \(2s^2 + s + 1 = 0\)

\begin{center}\rule{0.5\linewidth}{0.5pt}\end{center}

\end{document}
